% GENERATED FILE. Edit canonical lessons, then run npm run generate:books.
% canonical-source-hash: fnv1a64:c6299828762961c0
% canonical-lessons: ZH-C521-shua-scrub, ZH-C521-feng-sew, ZH-C521-buzhi, ZH-C521-zhengqi, ZH-C521-weisheng

\chapter{Keeping House}
\label{ch:zh-keeping-house}
\hlchaptermodality{521}

\begin{chapteropening}
\textbf{By the end of this chapter:} \emph{I can talk about scrubbing, mending, arranging a room, keeping it tidy and keeping a kitchen clean.}

Complete the last lesson of chapter 521: I can talk about scrubbing, mending, arranging a room, keeping it tidy and keeping a kitchen clean.
\end{chapteropening}

\section[shuā]{shuā — to scrub; to brush clean}

\label{lesson:ZH-C521-shua-scrub}



\begin{quote}
Before the new one: say the Mandarin for central heating, then the Mandarin for to use little electricity.
\end{quote}

\subsection*{You'll want to know: shuā}
\emph{shuā} — \textquotedblleft{}to scrub; to brush clean\textquotedblright{}.

\begin{quote}
\emph{wǒ shuā le yí ge xiǎoshí de guō} — I spent an hour scrubbing the pot.
\end{quote}

Its characters come later; for now it is a word for the ear.

\begin{cousinweb}
\emph{shuā} is inherited from Middle Chinese \emph{srwaet} (Guangyun, in Baxter's transcription), glossed 'to wipe clean'; the dictionary also lists \emph{srjwet}, glossed 'to sweep'. The level tone of shuā cannot be read off \emph{srwaet}: words that once ended in t, as this one did, are spread over all four Mandarin tones.
\end{cousinweb}

\begin{grammarlens}[title={What you've built}]
One more word for saying how it was.
\end{grammarlens}

\subsection*{Your turn}
\begin{itemize}
\raggedright
  \item \emph{Say it:} \emph{shuā}
  \item \emph{Say it:} \emph{shuā}, once more
  \item \emph{Say it:} \emph{wǒ shuā le yí ge xiǎoshí de guō} — I spent an hour scrubbing the pot.
  \item \emph{Recall:} say the Mandarin for central heating, then the Mandarin for to use little electricity, then say \emph{shuā} again
\end{itemize}

\begin{tcolorbox}[breakable,colback=teal!4,colframe=teal!35!black,title={Before you move on}]
What does shuā mean? (\textquotedblleft{}To scrub; to brush clean\textquotedblright{}.) Say it once more.
\end{tcolorbox}
\section[féng]{féng — to sew; to stitch}

\label{lesson:ZH-C521-feng-sew}



\begin{quote}
Before the new one: say the Mandarin for to use little electricity, then the Mandarin for to scrub clean.
\end{quote}

\subsection*{You'll want to know: féng}
\emph{féng} — \textquotedblleft{}to sew; to stitch\textquotedblright{}.

\begin{quote}
\emph{māma bāng wǒ féng kùzi féng le bàntiān} — Mum sewed up my trousers for me, and it took her ages.
\end{quote}

Its characters come later; for now it is a word for the ear.

\begin{cousinweb}
\emph{féng} is inherited from Middle Chinese \emph{bjowng} (Guangyun, in Baxter's transcription), glossed 'to stitch'; its departing-tone reading, \emph{bjowngH}, is glossed 'a seam'.
\end{cousinweb}

\begin{grammarlens}[title={What you've built}]
One more word for saying how it was.
\end{grammarlens}

\subsection*{Your turn}
\begin{itemize}
\raggedright
  \item \emph{Say it:} \emph{féng}
  \item \emph{Say it:} \emph{féng}, once more
  \item \emph{Say it:} \emph{māma bāng wǒ féng kùzi féng le bàntiān} — Mum sewed up my trousers for me, and it took her ages.
  \item \emph{Recall:} say the Mandarin for to use little electricity, then the Mandarin for to scrub clean, then say \emph{féng} again
\end{itemize}

\begin{tcolorbox}[breakable,colback=teal!4,colframe=teal!35!black,title={Before you move on}]
What does féng mean? (\textquotedblleft{}To sew; to stitch\textquotedblright{}.) Say it once more.
\end{tcolorbox}
\section[bùzhì]{bùzhì — to lay out a room; to decorate a room}

\label{lesson:ZH-C521-buzhi}



\begin{quote}
Before the new one: say the Mandarin for to scrub clean, then the Mandarin for to sew by hand.
\end{quote}

\subsection*{You'll want to know: bùzhì}
\emph{bùzhì} — \textquotedblleft{}to lay out a room; to decorate a room\textquotedblright{}.

\begin{quote}
\emph{xīn fángzi wǒmen bùzhì le liǎng tiān} — We spent two days arranging the new flat.
\end{quote}

Its characters come later; for now it is a word for the ear.

\begin{cousinweb}
\emph{bùzhì} is built from \emph{bù} 'to spread' and \emph{zhì} 'to place', both inherited from Middle Chinese, \emph{puH} and \emph{triH} (Guangyun, in Baxter's transcription), which glosses the first 'to spread everywhere' and the second 'to set in place'.
\end{cousinweb}

\begin{grammarlens}[title={What you've built}]
One more word for saying how it was.
\end{grammarlens}

\subsection*{Your turn}
\begin{itemize}
\raggedright
  \item \emph{Say it:} \emph{bùzhì}
  \item \emph{Say it:} \emph{bùzhì}, once more
  \item \emph{Say it:} \emph{xīn fángzi wǒmen bùzhì le liǎng tiān} — We spent two days arranging the new flat.
  \item \emph{Recall:} say the Mandarin for to scrub clean, then the Mandarin for to sew by hand, then say \emph{bùzhì} again
\end{itemize}

\begin{tcolorbox}[breakable,colback=teal!4,colframe=teal!35!black,title={Before you move on}]
What does bùzhì mean? (\textquotedblleft{}To lay out a room; to decorate a room\textquotedblright{}.) Say it once more.
\end{tcolorbox}
\section[zhěngqí]{zhěngqí — neat; orderly}

\label{lesson:ZH-C521-zhengqi}



\begin{quote}
Before the new one: say the Mandarin for to sew by hand, then the Mandarin for to lay out a room.
\end{quote}

\subsection*{You'll want to know: zhěngqí}
\emph{zhěngqí} — \textquotedblleft{}neat; orderly\textquotedblright{}.

\begin{quote}
\emph{tā de fángjiān bǐ wǒ de zhěngqí} — His room is tidier than mine.
\end{quote}

Its characters come later; for now it is a word for the ear.

\begin{cousinweb}
\emph{zhěngqí} is built from \emph{zhěng} and \emph{qí}, both inherited from Middle Chinese, \emph{tsyengX} and \emph{dzej} (Guangyun, in Baxter's transcription), which glosses each with the other: 'upright, even'.
\end{cousinweb}

\begin{grammarlens}[title={What you've built}]
One more word for saying how it was.
\end{grammarlens}

\subsection*{Your turn}
\begin{itemize}
\raggedright
  \item \emph{Say it:} \emph{zhěngqí}
  \item \emph{Say it:} \emph{zhěngqí}, once more
  \item \emph{Say it:} \emph{tā de fángjiān bǐ wǒ de zhěngqí} — His room is tidier than mine.
  \item \emph{Recall:} say the Mandarin for to sew by hand, then the Mandarin for to lay out a room, then say \emph{zhěngqí} again
\end{itemize}

\begin{tcolorbox}[breakable,colback=teal!4,colframe=teal!35!black,title={Before you move on}]
What does zhěngqí mean? (\textquotedblleft{}Neat; orderly\textquotedblright{}.) Say it once more.
\end{tcolorbox}
\section[wèishēng]{wèishēng — hygiene; cleanliness}

\label{lesson:ZH-C521-weisheng}



\begin{quote}
Before the new one: say the Mandarin for to lay out a room, then the Mandarin for neat and orderly.
\end{quote}

\subsection*{You'll want to know: wèishēng}
\emph{wèishēng} — \textquotedblleft{}hygiene; cleanliness\textquotedblright{}.

\begin{quote}
\emph{chúfáng yídìng yào zhùyì wèishēng} — A kitchen has to be kept really clean.
\end{quote}

Its characters come later; for now it is a word for the ear.

\begin{cousinweb}
\emph{wèishēng} is built from \emph{wèi} 'to guard' and \zh{生} 'life', both inherited from Middle Chinese, \emph{hjwejH} and \emph{srjaeng} (Guangyun, in Baxter's transcription), which glosses the first 'to protect' and the second 'to grow'. Guarding life as a word for hygiene is modern, and its path is not settled.
\end{cousinweb}

\begin{grammarlens}[title={What you've built}]
One more word for saying how it was.
\end{grammarlens}

\subsection*{Your turn}
\begin{itemize}
\raggedright
  \item \emph{Say it:} \emph{wèishēng}
  \item \emph{Say it:} \emph{wèishēng}, once more
  \item \emph{Say it:} \emph{chúfáng yídìng yào zhùyì wèishēng} — A kitchen has to be kept really clean.
  \item \emph{Recall:} say the Mandarin for to lay out a room, then the Mandarin for neat and orderly, then say \emph{wèishēng} again
\end{itemize}

\begin{tcolorbox}[breakable,colback=teal!4,colframe=teal!35!black,title={Before you move on}]
What does wèishēng mean? (\textquotedblleft{}Hygiene; cleanliness\textquotedblright{}.) Say it once more.
\end{tcolorbox}
